\documentclass[11pt]{article}
\usepackage[a4paper,margin=25mm]{geometry}
\usepackage{amsmath,amssymb,amsthm,booktabs,array}
\usepackage[T1]{fontenc}
\usepackage{lmodern}
\usepackage[colorlinks=true,linkcolor=blue,urlcolor=blue,citecolor=blue]{hyperref}
\usepackage{microtype}
\hypersetup{pdftitle={Structure, Exact Algorithms and Information Limits in a Unit-Latency Reorder-Buffer Model},pdfsubject={ILP1-20261014 research draft; originality unconfirmed}}
\newtheorem{theorem}{Theorem}
\newtheorem{proposition}[theorem]{Proposition}
\newtheorem{corollary}[theorem]{Corollary}
\newcommand{\T}{T^*}
\newcommand{\bs}{B_{\mathrm{sat}}}
\newcommand{\ws}{W_{\mathrm{sat}}}
\title{Structure, Exact Algorithms and Information Limits\\in a Unit-Latency Reorder-Buffer Model}
\author{}
\date{Research draft, 14 October 2026 (Asia/Singapore)}
\begin{document}
\maketitle
\vspace{-2em}
\begin{center}\small ILP1-20261014. Correctness and historical novelty are separate questions.\\
No author list, publication license, submission, or DOI is determined by this draft.
\end{center}
\begin{abstract}
We study offline and online unit-job scheduling with eager in-order admission,
bounded unretired occupancy, and in-order retirement. We prove an exact
execution-time reduction to precedence scheduling by adding every edge from
index $i$ to index $j$ with $i+B\le j$. This is a specialized formal bridge
to established resource-event graph methods, not a claim to have invented them.
We give exact algorithms and independently checked certificates, and exhibit
an exact polynomial optimizer for arbitrary issue width $W$ when the window
has capacity $B=W+1$. Perfect saturation then has a linear structural test.
More generally, a horizon with $S=HW-N$ unused issue slots confines every
feasible trace to a strip of width $2(B-W)+S$, yielding fixed-parameter
feasibility in $(B-W)+S$. We also exhibit
two-input programs whose minimum sufficient window falls from $N$ to $3$ when
issue width increases from $2$ to $3$. This unbounded threshold change costs
only one extra cycle at window two, highlighting the distinction between
exact and approximate saturation. A separate chain-prefix family has linear
window requirements even at fixed positive relative tolerance. Under a precise hidden-future information
model, a universally optimal online policy exists exactly when $W=1$ or
$W\ge B$. At $B=3,W=2$, a simple randomized policy has the sharp strict
ratio $5/4$, proved by a joint potential invariant. A two-input repeated
construction preserves this lower bound with an additive constant.
Lean 4.34.1 verifies the operational reduction, existence of an
optimum, an executable exact enumeration algorithm, resource caps, Bellman
certificates, the actual-trace strip bound, and a four-job arbitrary-policy online obstruction. General
resource families and online classifications have independent written review
but are not yet fully formalized. Historical originality remains unconfirmed.
\end{abstract}

\section{Introduction}
For a fixed program, increasing processor resources cannot worsen its optimal
duration. This does not imply that the least amount of one resource needed to
attain an optimum changes monotonically with another resource. Nor does the
existence of an offline optimum imply that a finite-information scheduler can
always select it. We separate these questions in a deliberately narrow model.

The main structural examples concern exact resource thresholds, rather than
real hardware speedups. A separate algorithmic result exploits how completed
but unretired jobs remove scheduling choices, yielding polynomial exact
optimization for a window one entry wider than issue width. Another result classifies the resource regimes in
which hidden future dependencies obstruct universal online optimality. The
foundational reduction and algorithms are included to make these claims
precise, reproducible, and amenable to formal checking. The result catalogue
and complete audit records accompany this draft in the research workstream.

\section{Formal model}
The program consists of $N$ fixed operations with indices $0,\ldots,N-1$.
Its true dependency graph $E$ has only forward edges $i<j$. Operations are
deterministic, have unit latency, and execute exactly once. Two-input,
one-output operations suffice for the principal examples. SSA or versioned
registers remove naming dependencies; this assumes pure operations and
preserves producer values. Instruction elimination and algebraic rewriting
are excluded; an isolated operation may also produce an observable output.

At a cycle boundary let $D$ be the set of completed operations and let $r$
be the least incomplete index, or $N$ at termination. With $B,W\ge1$, the
eager window is $[r,\min(N,r+B))$. Completed entries behind $r$ remain in
the window until earlier operations finish. In one cycle any subset of at
most $W$ ready, incomplete window entries may issue. A predecessor must have
completed in a previous cycle. Issued jobs complete at the cycle end; the
whole completed prefix retires and the next cycle refills immediately.
There is no same-cycle dependent bypass.

A schedule is a map $t:\{0,\ldots,N-1\}\to\mathbb N$ of issue times.
For nonempty programs its duration is $1+\max_j t(j)$; empty programs take
zero cycles and have no defined IPC. Let $\T(E,B,W)$ be the minimum duration
and $IPC^*=N/\T$. Both resources may be capped at $N$ for nonempty inputs;
effective width is at most $B$.

An offline planner sees the entire graph. An online planner sees only
admitted operations, their incoming true dependencies, and execution history.
It may know $N$, but receives no oracle for future uses of visible values.
All online lower bounds below use this observation model.

\section{Related work and claim boundaries}
Wall's report studies instruction-level parallelism under idealization
\cite{wall}; Dubey, Adams, and Flynn analyze instruction windows
\cite{dubey}. The latter full text was not obtained in this audit.
Tullsen and Calder already add window-edge dependencies to model ordered
window admission in a 1998 report \cite{tullsen}.
Fields, Rubin, and Bod\'ik use execution-event DAGs with resource edges
\cite{fields}; later work explicitly records FIFO/ROB commit-to-dispatch
constraints \cite{navada}. Our all-prefix execution-only reduction must be
read as a specialization of this prior modeling tradition.

The all-prefix constraints also occur as discrepancy precedences in
Cire--van Hoeve's sequencing MDD formulation, whose section9 gives an
exponential-in-discrepancy width bound \cite{cire}. Thus neither these
constraints nor a general finite-frontier state representation is new.
Our current algorithm questions instead concern variable $W$, the difference
$B-W$, and a total idle-slot budget. Recent work parameterizes scheduling by
time-window overlap or incomparable-graph degeneracy; the workstream audit
records these sources and access limits. Those parameters can grow with $W$
even when the parameters here are fixed, so their inspected bounds do not
directly imply the results below. This distinction is not novelty clearance.

Bellman recursion, finite exhaustive search, precedence longest paths and
Graham's list-scheduling guarantee are classical. At width two the reduced
problem is covered by Coffman--Graham's polynomial algorithm
\cite{cg,graham}. We do not claim a new solution of unrestricted precedence
scheduling. A September 2026 OpenAI artifact claims a polynomial
three-machine result with Lean sources \cite{oai}; its actual endpoint was
inspected but its full dependency closure was not independently compiled here.
Neither unconditional acceptance nor a statement that this question simply
remains open is justified by our current audit.

Hidden-successor adversaries and advice-counting methods also have prior art;
the workstream literature notes compare different release and observation
models. No exact equivalent of all the threshold formulas below was located,
but the search is incomplete. All prospective novelty is unconfirmed.

\section{Static reduction and exact algorithms}
\begin{theorem}[Exact window representation]
Define
\[
 E_B=E\cup\{(i,j):i+B\le j<N\}.
\]
A time assignment is legal in the ROB model if and only if it strictly
respects $E_B$ and every time layer has at most $W$ jobs.
\end{theorem}
\begin{proof}
When $j$ issues, every index $i\le j-B$ must already have retired, hence
completed in a previous cycle. This gives all the artificial precedences;
true data and issue-width conditions are necessary as well.

Conversely replay the time layers. Before $t(j)$ all indices $i\le j-B$
have completed. Thus the retirement frontier satisfies $r>j-B$ when
$j\ge B$, and admission is immediate when $j<B$. Since $j$ is still
incomplete, $r\le j<r+B$. It is in the window, all true predecessors are
ready, and the proposed layer fits the issue width. Induction over time
proves legality, including schedules with idle cycles.
\end{proof}
Only adding the offset edge $(j-B,j)$ is insufficient: for $N=4,B=2$ and
$E=\varnothing$, the times $(1,0,2,1)$ satisfy both offset edges but violate
the required edge $0\to3$.

\begin{proposition}[Exact search and certificates]
An optimal finite schedule exists. If $A(D)$ is the ready window set, its
value satisfies
\[
 V(D)=1+\min_{\varnothing\ne S\subseteq A(D),\ |S|\le W}V(D\cup S),
 \qquad V(\{0,\ldots,N-1\})=0.
\]
A legal $K$-cycle witness plus a potential $v$ on every reachable state with
$v(\varnothing)=K$, $v(\mathrm{all})=0$ and
$v(D)\le1+v(D\cup S)$ for every legal transition proves optimality.
\end{proposition}
\begin{proof}
The earliest incomplete operation is ready because all its true predecessors
have smaller indices. Issuing it repeatedly finishes within $N$ cycles.
Idle steps are unnecessary, and each nonempty step strictly increases $|D|$.
Induction on $N-|D|$ proves the recurrence and constructs its attaining
schedule. Telescoping the potential inequality along any execution gives
the lower bound $K$; the legal witness gives the opposite bound.
\end{proof}
The independent checker enumerates successors itself and rejects omitted
reachable states. This is a full-state certificate, not a compact one.
There are at most $1+N2^{B-1}$ completed-set states. Labeling each frontier
position completed, selected, or unselected bounds aggregate transitions
by $O(N3^B)$ before polynomial feasibility-checking factors.

For $W\ge B$, any $E_B$ antichain lies in fewer than $B$ index gaps and has
at most $B$ vertices. Earliest precedence times are therefore feasible and
$\T(E,B,W)=h(E_B)$, where height counts vertices. The recurrence
\[
 e(j)=1+\max\left(\max_{(i,j)\in E}e(i),\ \max_{i\le j-B}e(i)\right),
\]
with empty maximum $-1$, is evaluated in $O(N+|E|)$ using prefix maxima.
For $W\ge N$, binary search on $B$ computes $\bs$ in
$O((N+|E|)\log N)$. A height-optimal schedule, a true critical path, and a
longer augmented path at $B-1$ (only when $B>1$) form an $O(N)$ certificate
of its minimality. At $B=1$ no smaller positive window needs exclusion.

For any fixed true-dependency-respecting schedule, define
\[
 \beta(t)=\max\bigl(\{1\}\cup
 \{j-i+1:i<j,\ t(i)\ge t(j)\}\bigr),\quad
 \omega(t)=\max_c|t^{-1}(c)|.
\]
The reduction gives legality exactly when $B\ge\beta(t)$ and
$W\ge\omega(t)$. Minimizing these footprints over the appropriate optimal
schedules characterizes saturation, but is not itself an efficient general
algorithm.

\section{Polynomial optimization with one extra window entry}
Call an action non-delay if it issues $\min(W,|R(D)|)$ jobs, where $R(D)$
is the ready set. This specifies a batch size, not a priority order.

\begin{proposition}[Completion dominance and non-delay optimality]
If two reachable states have completed sets $D\subseteq D'$, every
continuation from $D$ can be simulated from $D'$ in no more cycles.
Consequently every reachable state has a non-delay optimal continuation.
\end{proposition}
\begin{proof}
Replay the original actions, deleting jobs already complete in the simulating
state. True predecessors stay complete. Retirement advances from $r$ to
$r'\ge r$; a retained job satisfies $j\ge r'$ because it is incomplete, and
$j<r+B\le r'+B$ because it was originally admitted. Thus it remains admitted
and ready. Deletion cannot exceed issue width. This preserves completed-set
inclusion at each boundary and finishes no later. Enlarge an optimal first
action to the non-delay size and apply this simulation. Induction on the
number of unfinished jobs gives a non-delay optimum at every later state.
\end{proof}

Write $h(D)=|D|-r(D)$ for completed but unretired jobs. They still occupy
window entries, so $|R(D)|\le B-h(D)$. If $B=W+1$ and $h\ge1$, the
non-delay action is therefore uniquely determined: issue every ready job.
Choices occur only at clear-prefix states $D=[0,i)$, and only when all
$W+1$ window entries are ready.

\begin{theorem}[Exact prefix dynamic program]
For arbitrary finite forward dependencies and arbitrary input $W\ge1$,
the optimum at $B=W+1$ and an attaining schedule can be computed in
\[
 O\bigl(N\min(B,N)(N+|E|)\bigr)\text{ time},\qquad O(N+|E|)\text{ space}
\]
using indexed-graph operations. No divisibility assumption on $N$ is needed.
\end{theorem}
\begin{proof}
At prefix $i$, enumerate its non-delay first actions: one if at most $W$
jobs are ready, otherwise the $B$ choices of an omitted ready job.
After each action, simulate the unique non-delay actions until the first
new clear prefix $j$. Every nonterminal cycle issues at least the oldest
job when its action is unique; the first action is also nonempty. Hence
this macro ends in at most $N-i$ cycles, has length $\ell$, and $j>i$.
The non-delay optimum must follow one such macro. Conversely every macro
is a legal continuation. Thus
\[
 F_N=0,\qquad F_i=\min_{\text{first actions }A}
       \bigl(\ell(i,A)+F_{j(i,A)}\bigr).
\]
Backward induction gives the exact remaining optimum at every prefix.

There are at most $N\min(B,N)$ macros. For each, initialize predecessor
counters from its starting prefix in $O(N+|E|)$ time. Monotone admission
and retirement pointers admit and complete each job once during that macro;
each outgoing edge decrements a counter once. Newly ready jobs wait until
the next cycle. Thus a macro costs $O(N+|E|)$. Store only its winning first
omission, cost and endpoint at each prefix, then replay winning macros for
the schedule. This gives the bounds. Empty programs return zero, and
arbitrarily large $W$ does not cause iteration over absent jobs.
\end{proof}

This is an offline algorithm: evaluating choices uses dependencies that may
be outside the current execution window. A forced macro can have length
linear in $N$ even at fixed $B=4,W=3$, so a constant macro-length assumption
would be invalid.

For general normalized excess $c=B-W\ge1$, branching requires $h<c$.
There are at most $N\sum_{h=0}^{c-1}\binom{B-1}{h}$ branching states.
The total number of outgoing macros is at most
\[
 N\sum_{h=0}^{c-1}\binom{B-1}{h}\binom{B-h}{W}
 \le N(2^c-1)\binom{B}{c},
\]
by replacing $\binom{B-1}{h}$ with $\binom{B}{h}$ and applying the identity
$\binom{B}{h}\binom{B-h}{c-h}=\binom{B}{c}\binom{c}{h}$.
Compress each unique-action segment, including an initial forced segment
if necessary. Completed-set size strictly increases along every macro,
giving an exact acyclic state graph. This is an XP algorithm in $c$,
not a fixed-parameter algorithm in $c$ alone. At $c=0$ every action is unique.

\subsection{A linear perfect-saturation criterion}
Suppose $W\ge2$, $B=W+1$ and $N=kW$. Partition the jobs into $k$
consecutive blocks of $W$ entries. A singleton macro executes one block
in one cycle. A pair macro swaps the last job of its left block with the
first job of its right block, then executes the resulting sets in two
successive cycles. A macro is legal when every internal true edge points
strictly from an earlier to a later macro cycle.

\begin{theorem}[Exact-work-bound tiling]
$\T=k$ if and only if these blocks can be tiled by legal singleton and
pair macros. Every $k$-cycle schedule has this decomposition. The decision
and a positive or negative certificate take $O(N+|E|)$ graph operations.
\end{theorem}
\begin{proof}
Every cycle must issue $W$ jobs. At a clear-prefix state with at least two
blocks remaining, a full action omits exactly one of $W+1$ entries, at
one-based position $r$. Only the first $r-1$ entries retire, leaving at
most $r$ unfinished window entries next cycle. Full width forces $r\ge W$.
For $r=W+1$ this is a singleton macro. For $r=W$ the next full action is
forced and gives exactly the pair macro. Both end at a clear prefix.
The final single block, if present, is a singleton. Conversely legal macros
satisfy admission and return to clear prefixes, so a tiling gives $k$ cycles.

Make a path on blocks, keeping an edge iff its pair macro is legal and
marking a vertex iff its singleton is legal. An even component can be paired.
An odd component can be covered iff it has a marked odd local position:
use that singleton and pair both sides; conversely the leftmost singleton
of any tiling must have odd position. Each true edge invalidates at most
one singleton and two adjacent-pair macros. One graph scan and one path
scan therefore suffice. For rejection, supply an odd interval isolated by
illegal boundary pair macros and with all odd-position singletons illegal.
One actual true edge witnesses each illegality; internal pair flags need
not be checked. No full-state certificate is needed.
\end{proof}

\section{A sharp completion corridor near the work bound}
Let $1\le W\le B$, choose a horizon $H$ with $WH\ge N$, and put
$c=B-W$, $S=WH-N$. A trace finishing earlier is padded with terminal idle
cycles, so its total unused issue capacity is exactly $S$.

\begin{theorem}[Actual-trace corridor]
At every boundary $0\le t\le H$, every such execution satisfies
\[
 [0,\max(0,tW-c-S))\subseteq D_t\subseteq[0,\min(N,tW+c)),
 \qquad tW-S\le |D_t|\le tW.
\]
\end{theorem}
\begin{proof}
For $t<H$, let $r$ be the first incomplete index at $t$. All jobs completed
by $t+1$ lie below $r+B$: earlier admissions had no larger frontier, and
the next action uses the current window. At most $(H-t-1)W$ jobs fit in
the remaining cycles, hence $N\le r+B+(H-t-1)W$ and $r\ge tW-c-S$.
At $t=H$, use $r=N$ directly. All indices below $r$ are complete.

A job completed in cycle $s\le t$ was admitted below
$r_{s-1}+B\le(s-1)W+B=sW+c\le tW+c$.
Finally $|D_t|\le tW$ and $N-|D_t|\le(H-t)W$ give the count interval.
The argument allows idle cycles, truncated windows and empty programs.
\end{proof}

This theorem, including counts derived from actual cycle/lane injections,
is verified in Lean for actual ROB traces. The following combinatorial
count and algorithm currently have reviewed written proofs.

\begin{theorem}[Sharp candidate count and parameterized feasibility]
Each layer has at most
\[
 M(c,S)=\sum_{j=c}^{c+S}\binom{2c+S}{j}\le2^{2c+S}
\]
candidate completed sets. This bound is sharp. Exact horizon feasibility
and an attaining schedule can be computed in
$2^{O(c+S)}\operatorname{poly}(N+|E|)$ time.
\end{theorem}
\begin{proof}
Embed the variable strip into the untruncated interval
$[tW-c-S,tW+c)$. Treat negative indices as fixed complete, and indices
at least $N$ as fixed incomplete. If $d=tW-|D_t|$, its selected subset
has size $c+S-d$, where $0\le d\le S$. This is an injection into the
families counted by $M$; boundary truncation only removes candidates.

For sharpness choose $W\ge\max(1,c+S)$, $H=2$, $N=2W-S$ and independent
jobs. Let $L=W-c-S$, $U=W+c$. For every subset $A\subseteq[L,U)$ of
size between $c$ and $c+S$, issue $D=[0,L)\cup A$ in cycle one.
It is initially admitted and has size between $W-S$ and $W$.
Its retired prefix is at least $L$, so its next frontier is at least
$L+B=N$. At most $W$ jobs remain, all independent, and finish in cycle two.
All $M$ distinct middle states therefore occur on one input. This is not
a computational lower bound; the input itself is easy.

Generate the layer candidates and test every consecutive-layer pair for
the actual ROB transition: containment of completed sets, at most $W$ new
jobs, admission using the actual first missing index, and all true parents
already complete. Empty additions are allowed. Every edge is sound; the
corridor puts every feasible trace inside the graph. Reachability and
predecessors decide feasibility and recover a schedule. There are at most
$M^2$ pairs per layer; graph/mask arithmetic is polynomial. If $H\ge N$,
use the sequential witness instead, so a binary-encoded horizon is harmless.
No step enumerates all $2^B$ ready subsets.
\end{proof}

At $c=1,S=0$ there are only two candidate forms: a clear prefix or its
boundary swap. A swapped state cannot be followed by another swapped state
at full width because the next required final index is outside admission.
This recovers the preceding tiling result. For $S>0$, appending independent
dummy jobs is not a reduction to $S=0$: at $B=4,W=3$, five jobs with sole
true edge $0\to1$ finish in two cycles, but adding a sixth independent job
makes two cycles impossible.

\section{Unbounded changes in exact resource thresholds}
For nonempty programs define
\[
 \bs(E,W)=\min\{B:\T(E,B,W)=\T(E,N,W)\},
\]
\[
 \ws(E,B)=\min\{W:\T(E,B,W)=\T(E,B,B)\}.
\]
Optimal duration is nonincreasing in each resource, since more window space
removes artificial edges and more issue capacity relaxes layer size. The
threshold functions do not inherit that monotonicity.

\begin{theorem}[A fixed-width window collapse]
For every integer $K\ge2$ there is a program $P_K$ of $N=2K$ operations,
with true indegree and outdegree at most two, such that
\[
 \bs(P_K,2)=2K,\qquad \bs(P_K,3)=3.
\]
Both unrestricted-window optima equal $K$.
\end{theorem}
\begin{proof}
Job $0$ is isolated. The other jobs are pair layers
$\{1,2\},\{3,4\},\ldots,\{2K-3,2K-2\}$ followed by singleton
$\{2K-1\}$. Between successive layers add all forward edges. Every
nonisolated job lies on a $K$-vertex path. Consequently any $K$-cycle
schedule forces the job in layer $\ell$ into cycle $\ell-1$.

At width two every pair layer fills its cycle; job $0$ must share the final
cycle with job $2K-1$. This is legal at $B=N$ and meets the height bound.
At every $B<N$, however, the artificial edge $0\to2K-1$ prohibits that
common time. Hence $\bs(P_K,2)=N$.

At $B=W=3$, issue $\{0,1,2\}$ first, then successive pairs, then the final
singleton. Every batch is admitted and the duration is $K$. At $B\le2$,
job $2$ is not admitted initially, contradicting its forced initial time
in every $K$-cycle schedule. Thus the second threshold is exactly three.
\end{proof}

The additive threshold drop $2K-3$ and the ratio $2K/3$ are unbounded, even
for fixed issue widths. This does \emph{not} imply an unbounded performance
gain. At width two and every $2\le B<2K$, the exact duration is only $K+1$:
issue $\{0,1\}$, then $\{2\}$, then the remaining pair layers and the final
singleton. Window two already suffices for this schedule, and the theorem
excludes $K$ cycles. At window one the duration is $2K$.

Define $B_\varepsilon(E,W)$ by allowing duration at most
$(1+\varepsilon)\T(E,N,W)$. The exact curves for this family are
\[
\begin{array}{c|ccc}
 &0\le\varepsilon<1/K&1/K\le\varepsilon<1&\varepsilon\ge1\\\hline
 W=2&2K&2&1\\
 W=3&3&2&1
\end{array}
\]
Thus this example's huge exact threshold is fragile under fixed relative
tolerance as $K$ grows. The following separate construction shows that such
fragility is not universal.

\begin{theorem}[A depth-two width-saturation curve]
Let $m\ge2$ and $1\le r\le\lfloor m/2\rfloor$. Take $r$ initial isolated
roots, then $m$ required roots, then $m-r$ sinks. Partition the required
roots into $r$ pairs and $m-2r$ singletons, assigning each group to a
different sink. For $1\le B\le2m$,
\[
 \T(E,B,W)\le2\quad\Longleftrightarrow\quad
 B\ge m+r\ \text{ and }\ W\ge m+\min(r,2m-B).
\]
In particular, $\ws(E,B)=m+\min(r,2m-B)$ for $m+r\le B\le2m$.
\end{theorem}
\begin{proof}
All $m$ required roots must issue first, so the last one must be initially
admitted: $B\ge m+r$. Admission of the final job in the next cycle forces
the first $2m-B$ jobs to complete initially. The union of these jobs and
the required roots has size $m+\min(r,2m-B)$, giving the width bound.
Conversely issue exactly that union first. Its prefix retirement admits
the whole remaining program, every sink is ready, and at most $m$ jobs
remain for the second cycle. Any true edge excludes a one-cycle schedule.
\end{proof}
As $B$ grows from $2m-r$ to $2m$, the sufficient width decreases by $r$.
Original indegree is at most two and outdegree at most one.

\begin{theorem}[Robust linear window demand]
Let $Q_K$ consist of a chain in indices $0,\ldots,K-1$, followed by $K$
independent operations, where $K\ge1$. For all $B\ge1,W\ge2$,
\[
 \T(Q_K,B,W)=K+\left\lceil
 \frac{\max(0,K-B+1)}{\min(B,W)}\right\rceil.
\]
Writing $q=\lfloor\varepsilon K\rfloor$ for $\varepsilon\ge0$, its exact
approximate saturation threshold is
\[
 B_\varepsilon(Q_K,W)=\max\left(
 \left\lceil\frac{K+1}{q+1}\right\rceil,\ K+1-qW\right).
\]
For every fixed $W\ge2$ and $0\le\varepsilon<1/W$, this threshold is
$\Omega(N)$, although true indegree and outdegree are at most one.
\end{theorem}
\begin{proof}
The final chain operation completes no earlier than the end of the $K$th
cycle. At the beginning of its issue cycle, the retirement frontier is at
most $K-1$. Hence at most $B-1$ independent operations have even been
admitted, including any issued alongside this last chain operation. At least
$Q=\max(0,K-B+1)$ independent operations remain afterward, requiring
$\lceil Q/\min(B,W)\rceil$ further cycles. Delaying the chain cannot evade
this admission bound.

For attainment, run the chain without gaps. In each of its last
$\min(K,B-1)$ cycles also issue the next independent operation. These
operations are admitted as the chain prefix retires, and width at least two
suffices. After the chain finishes, the remaining $Q$ independent operations
can be issued in consecutive batches of size $\min(B,W)$. This proves the
duration formula, including $B=1$ and $B>N$.

At unrestricted window the optimum is $K$, so the allowed extra cycles are
$q$. The condition for a window $B$ is exactly
$K+1\le B+q\min(B,W)$. The identity
\[
 B+q\min(B,W)=\min((q+1)B,\ B+qW)
\]
turns this into the two lower bounds in the formula. In particular
$B_\varepsilon\ge K+1-W\lfloor\varepsilon K\rfloor
\ge(1-W\varepsilon)K+1$, proving the final claim.
\end{proof}
For width two and $0\le\varepsilon<1$, this simplifies to
$\max(2,K+1-2\lfloor\varepsilon K\rfloor)$; at $\varepsilon\ge1$ it is
one. At width one, the unrestricted baseline is instead $2K$, and the
relative saturation window is one for every nonnegative tolerance.

\begin{corollary}[A robust single-slot trade-off]
Fix an integer $W\ge2$ and $1/(W+1)\le\varepsilon<1/W$. On the same
family, $B_\varepsilon(Q_K,W)=\Omega(N)$ whereas
$B_\varepsilon(Q_K,W+1)\le W+1$. Thus increasing issue width by one can
reduce the approximate sufficient window by an unbounded factor at fixed
positive tolerance.
\end{corollary}
\begin{proof}
The lower bound follows from the theorem. Since
$q\ge\lfloor K/(W+1)\rfloor$, integer division gives
$K+1\le(W+1)(q+1)$. Both terms in the formula for width $W+1$ are
therefore at most $W+1$.
\end{proof}
For width two this applies to every fixed $1/3\le\varepsilon<1/2$.
This concerns approximate thresholds: the exact threshold of $Q_K$ is
$K+1$ for all $W\ge2$. Program order matters. Interleaving the same chain
with its independent operations allows $B=W=2$ to finish in $K$ cycles,
but reordering is excluded here. The relevant object is a DAG with its fixed
topological program order, not an unlabeled partial order alone.

\section{Information limits for online scheduling}
\begin{theorem}[Universal-optimality classification]
For fixed positive $B,W$, a deterministic online policy optimal on every
finite R1 program exists if and only if $W=1$ or $W\ge B$. The same
classification holds for randomized policies required to be optimal in
expectation on every input.
\end{theorem}
\begin{proof}
At $W=1$, repeatedly issue the oldest ready operation: every nonidle
execution takes $N$ cycles. At $W\ge B$, issue every ready operation.
The static reduction identifies this with augmented earliest-time layering,
which is optimal and needs no future information.

Suppose $B>W\ge2$ and put $L=\lceil(B-W)/(W-1)\rceil$. The first $B$
operations are roots. An $L$-job suffix chain starts from a hidden root
$x\in\{1,\ldots,B-1\}$. These inputs have identical initial views and
the same known length $B+L$.

For every $x$, issue any $W$ roots containing $0,x$ first. Then issue
each successive chain job along with up to $W-1$ smallest remaining roots.
Before suffix job $k$, the completed root prefix has length at least
$\min(B,1+(k-1)(W-1))\ge k$, so admission is legal. The capacity inequality
$W+L(W-1)\ge B$ finishes all roots. This realizes $L+1$ cycles and meets
the true chain lower bound.

Every deterministic first action either omits $0$, delaying suffix admission,
or includes $0$ and omits some possible $x$. The corresponding fixed input
cannot start its suffix in cycle one, and needs at least $L+2$ cycles.
For a randomized first action, average over uniformly chosen $x$. An action
containing $0$ hits at most $W-1$ of the $B-1$ choices, and one omitting
$0$ fails every choice. Positive average failure implies some fixed input
with expected duration strictly greater than its optimum.
\end{proof}

The lower-bound family uses only a true chain and isolated vertices. It is
an information obstruction, even without a computational complexity bound
on the policy. On this family, the exact randomized minimax ratio is
$1+(B-W)/((B-1)(L+1))$; choosing $0$ plus a uniform $(W-1)$-subset and
using one extra root-only cycle after a miss attains it. The complete
continuation and advice-counting proofs are in the accompanying theory note.

At $B=3,W=2$, the four-job choices $E=\{1\to3\}$ and
$E=\{2\to3\}$ force distinct optimal initial batches $\{0,1\}$ and
$\{0,2\}$. Every deterministic policy has ratio at least $3/2$ on one.
Graham's $2-1/W$ bound applies to every non-delay policy on the augmented
DAG, proving the global strict deterministic minimax ratio is exactly
$3/2$ at these resources. The next section sharpens the global randomized
value to exactly $5/4$. The finite witnesses establish strict ratios with
\emph{no additive constant}. A separate repeating construction below
also settles the definition allowing a fixed additive constant.

\section{The sharp randomized bound at a three-entry window}
Fix $B=3,W=2$. Policy R always issues the oldest unfinished operation.
When all three window entries are unfinished and ready, it chooses one of
the two younger operations uniformly, using fresh randomness. Otherwise it
issues all ready jobs. The following bound applies to every fixed program;
it does not give the adversary access to realized random choices.

\begin{theorem}[Sharp randomized guarantee]
Let $F$ be the offline optimum, and $E$ the expected duration of R. Then
\[
 4E\le N+3F\le5F.
\]
For nonempty inputs the global strict randomized online minimax ratio at
these resources is exactly $5/4$ against an oblivious adversary.
\end{theorem}
We give the complete upper-bound argument. The preceding four-job pair
supplies the lower bound: under the uniform distribution over its two
fixed programs, any initial action succeeds in two cycles with probability
at most $1/2$. Thus some fixed input has expected duration at least $5/2$
and optimum two.

First, when $B=W+1$, some optimal non-delay schedule includes the oldest
unfinished job every cycle. Choose the lexicographically least time vector
among all optimal schedules. It is non-delay because a ready job can move
into any unused issue slot. If oldest $a$ is omitted, all $W$ other window
entries must issue; retirement stalls and the next cycle issues only $a$.
Swap $a$ with any job of the first batch. The two-cycle endpoint is unchanged
and the time vector decreases, a contradiction. The swapped schedule need
not itself be non-delay: the minimum was over all optimal schedules.

For a completed-prefix state beginning at index $i$, let $F_i,E_i$ be its
optimal and expected remaining times, and $n_i=\max(0,N-i)$. Beyond the
program all three quantities are zero. Write $a=i,b=i+1,c=i+2,d=i+3$,
omitting nonexistent jobs. If just $a$ is ready, advance one job in one
cycle. If only younger job $b$ is ready, advance two jobs in one cycle.
Issuing $a,c$ forces a second cycle with $b$ and, if ready, $d$. This
advances $\ell=4$ jobs if $d$ exists and $b\to d$ is absent, otherwise
$\ell=3$. When both younger jobs are ready, optimality takes the minimum
of the two macro-actions $(1,i+2)$ and $(2,i+\ell)$; R takes their average.
The normalization just proved establishes completeness of this recurrence.

We also have $0\le F_i-F_{i+1}\le1$: completing an additional prefix job
cannot harm a continuation, while issuing oldest alone costs one cycle.
Define
\[
 K_i=3F_i+n_i-4E_i.
\]
We prove jointly, by descending induction, that $K_i\ge0$, and that
\[
 i<N,\quad F_i=F_{i+1}\quad\Longrightarrow\quad K_i+K_{i+1}\ge1.
 \tag{*}
\]
The restriction $i<N$ excludes two empty suffixes. At the final job,
$F_i=E_i=n_i=1$, hence $K_i=0$ and the equality antecedent is false.

A forced macro-action advancing $\ell$ jobs in $z$ cycles gives
$K_i=\ell-z+K_{i+\ell}$. The increments for $(\ell,z)=(1,1),(2,1),
(3,2),(4,2)$ are $0,1,1,2$. In a randomized state put
$p=i+2,q=i+\ell,x=F_p-F_q$. Direct substitution gives
\[
\begin{array}{c|c|c}
\ell&x&K_i\\\hline
3&0&(K_p+K_q)/2-1/2\\
3&1&(K_p+K_q)/2+1\\
4&0\text{ or }2&(K_p+K_q)/2\\
4&1&(K_p+K_q)/2+3/2
\end{array}
\]
These exhaust the possibilities because each consecutive $F$ difference
is zero or one. Every row is nonnegative by the tail induction hypothesis,
except possibly $\ell=3,x=0$. There $p$ is nonempty, $q=p+1$ and
$F_p=F_q$, so tail invariant (*) cancels the negative half. This proves
$K_i\ge0$ before proving (*) at the same index.

For (*), a forced singleton would give $F_i=1+F_{i+1}$, excluding its
antecedent. All other forced cases and the randomized cases $x=1$ give
$K_i\ge1$. It remains to handle three randomized cases, with $a,b,c$
all present and ready, hence with no true edges among them.

If $\ell=3,x=0$, then $F_c=F_d$. Job $d$ must exist, since $c$ is nonempty
and has positive value. The edge $b\to d$ therefore exists. At the suffix
beginning at $b$, only younger job $c$ is ready, so
$K_b=1+K_d\ge1$. Together with $K_a\ge0$ this proves (*).

If $\ell=4,x=2$, put $e=i+4$. Then $F_c=F_d+1=F_e+2$. Without the edge
$c\to d$, the legal first pair $c,d$ would give $F_c\le1+F_e$, a
contradiction. Thus $c\to d$ exists, and at suffix $b$ only younger job
$c$ is ready. Again $K_b=1+K_d\ge1$.

Finally suppose $\ell=4,x=0$, so $F_c=F_e$. If $c\to d$ exists the
preceding argument applies. Otherwise the pair $c,d$ is legal but costs
$1+F_e>F_c$. The recurrence forces the alternative $c,e$ to exist, be
ready and be optimal. In particular $c\to e$ and $d\to e$ are absent.
Its two-cycle target is $u\in\{e+1,e+2\}$ with $2+F_u=F_e$.
The one-step difference bound excludes $u=e+1$, so
\[
 F_e-F_{e+1}=F_{e+1}-F_{e+2}=1,\qquad F_d=F_e.
\]
At suffix $b$, all $b,c,d$ are ready. Its second macro-action $b,d$
also issues $e$ next cycle because $c\to e$ is absent. Its targets are
$d,e+1$, whose optimal values differ by one. This is the good
$\ell=4,x=1$ row, so $K_b\ge3/2$. This completes (*) and the induction.

At the initial state, $K_0\ge0$ gives $4E\le N+3F$. Width two gives
$N\le2F$, finishing the upper bound. On each member of the four-job pair,
R has $E=5/2,F=2$, attaining equality. Equivalently, R achieves in
expectation at least three quarters of the optimum number of double-issue
cycles. This pairing interpretation does not by itself establish an
equivalence to an arbitrary online matching model.

The general statement above has written proof and independent review.
Lean verifies its conditional joint-induction kernel and each local
min/mean potential identity. The remaining interface contains nine local
macro classifications and their structural successor equations; these
are not assumed nonnegativity conclusions. The connection from all ROB
programs to this interface, and the probability/scale semantics, are not
yet formalized. A prior exhaustive search of 2,796,203 effective local-edge
graphs through $N=12$ found no counterexample; it is supporting evidence,
not the proof. Historical novelty of the exact bound remains unconfirmed.

\section{Repeating the obstruction without separator cycles}
\begin{theorem}[Lower bounds with additive constants]
At $B=3,W=2$, the best randomized competitive factor against oblivious
inputs is $5/4$, even in the definition
$\mathbb E[T(P)]\le c\T(P)+\alpha$ with a fixed finite additive constant
$\alpha$. The deterministic factor in the same definition is $3/2$.
The lower-bound inputs have true indegree at most two.
\end{theorem}
\begin{proof}
For $q\ge1$, concatenate blocks
$(a_k,b_k,c_k,d_k)=(4k,4k+1,4k+2,4k+3)$, $0\le k<q$.
Within each block $d_k$ depends on a hidden choice $x_k\in\{b_k,c_k\}$.
For $k>0$, each of $a_k,b_k,c_k$ has exactly the two true predecessors
$c_{k-1},d_{k-1}$. All indices remain forward and true indegrees are at
most two.

The first job of a new block has artificial predecessors $a_{k-1},b_{k-1}$
and true predecessors $c_{k-1},d_{k-1}$. For the next job, window edges
already cover the first three previous jobs and its true edge covers the
fourth; for the third job, window edges cover all four. Thus no job of
the new block can issue before all jobs of the previous block complete.
This is strict cycle separation without a separator job or extra cycle.
Every block needs at least two cycles by width, and the batches
$\{a_k,x_k\}$ followed by the other root and $d_k$ attain two.
Consequently the offline optimum is $2q$.

Choose the hidden bits independently and uniformly before execution.
Before the first block action, $d_k$ is unadmitted because $a_k$ has not
retired. Earlier visible operands reveal no information about $x_k$;
future roots refer only to the fixed identities $c_k,d_k$. Conditional
on every previous observation and all policy random state, the current
bit is therefore still fair. Completion in two cycles forces one of the
two distinct first actions $\{a_k,b_k\}$ or $\{a_k,c_k\}$, according to
the hidden bit. Success probability is at most one half; every failure,
including idle or underfilled first choices, needs at least three cycles.
Separation and linearity of expectation give average duration at least
$5q/2$. For every $q$ there is thus a fixed input with expected duration
at least $5q/2$ and optimum $2q$. The inequality
$5q/2\le2cq+\alpha$ for every $q$ forces $c\ge5/4$.

For a deterministic policy, select each still-hidden bit to disagree
with its first action. This recursively determines one fixed complete
program, since the action cannot yet depend on that bit. Every block
needs at least three cycles, proving $3q\le2cq+\alpha$ and $c\ge3/2$.
This argument is only for deterministic policies; the randomized proof
used the oblivious product distribution. The preceding strict upper
bounds complete both equalities.
\end{proof}
The product-distribution and direct-sum methods are classical. Historical
priority of the exact two-input ROB composition is unconfirmed. The
general composition has two independent written reviews. Lean verifies
its true indegree bound, arbitrary-block separation, exact offline value
$2q$, and hidden-job admission barrier in the actual ROB model, including
$q=0$. Conditioning on observation histories and the probabilistic lower
bound are still written arguments.

\section{Larger windows and a failed policy extension}
\begin{proposition}[A hidden-subset family]
Let $W\ge2$, $1\le m\le W-1$, $W<B\le2W-m$, and let $d\ge1$ satisfy
$W-m\le dm$. There is a family of
$M=\binom{B-m}{W-m}$ programs with $N=B+m$, true indegree at most $d$,
outdegree at most one and depth at most two, on which the exact randomized
minimax ratio is $3/2-1/(2M)$. Exactly $\lceil\log_2M\rceil$ bits of
fixed-length pre-decision advice are necessary and sufficient for a
deterministic policy always to be optimal within this promised family.
\end{proposition}
\begin{proof}
The first $B$ jobs are roots. Put $F=\{0,\ldots,m-1\}$ and
$U=\{m,\ldots,B-1\}$. Choose a hidden subset $A\subset U$ of size $W-m$
and partition it among the $m$ suffix jobs, each with at most $d$ parents.
Empty parts are allowed. All initial observations agree. Any two-cycle
completion must issue $F$ first to admit the last job, and $A$ first to
satisfy every suffix dependency. Thus its first action is uniquely
$F\cup A$, filling width $W$. Conversely the remaining $B-W$ roots and
$m$ suffix jobs fit cycle two. The optimum is two.

The $M$ distinct required actions give success probability at most $1/M$
under uniform hidden input, and therefore expected duration at least
$3-1/M$ on some fixed input. Guessing $A'$ uniformly attains this value:
issue $F\cup A'$, then every remaining root and ready suffix job, which
fit because $B-W+m\le W$. On a wrong guess the blocked sinks all fit
cycle three. Advice needs $M$ messages to distinguish the first actions,
and encoding $A$ suffices. This is the classical distinguishability
counting argument applied to this family.
\end{proof}
The attaining policy is promised the family; this is a global-model lower
bound, not a global upper bound. At $B=W+1$, $M=W+1-m$.
Taking $W=2q,m=q,d=1$ gives a disjoint-edge true graph and lower bound
$3/2-1/(2(q+1))$, tending to $3/2$ as the resources grow. This is distinct
from an asymptotic input-length claim at fixed resources.

At $B=4,W=3$, uniform omission of one younger job when all four are ready
does not guarantee $4/3$. For $N=9$ with zero-based true edges
\[
 \{1\to4,\ 2\to4,\ 2\to5,\ 5\to8\},
\]
the optimum is three, attained by
$\{0,1,2\};\{3,4,5\};\{6,7,8\}$. The uniform policy has durations
three, four and five with probabilities $1/9,2/3,2/9$, respectively.
Its expected ratio is $37/27>4/3$. An independent constraint solver and
direct probability-tree replay verify these exact values. Exhaustive
effective local-edge search through $N=8$ found no smaller counterexample
to this policy. A different online policy might still attain $4/3$;
that minimax question remains open here.

\section{Computational and formal verification}
The operational Python solver uses completed-set memoization. A separately
implemented time-assignment constraint solver shares neither that recurrence
nor its state encoding. Operational replay and Bellman checking are also
separate implementations. Evidence preserves seeds, source hashes, timings,
and exact tested ranges. The following datasets overlap and must not be
added as if they were disjoint experiments.

\begin{center}\small
\begin{tabular}{p{0.29\linewidth}p{0.62\linewidth}}
\toprule
Audit & Actual scope\\\midrule
Root exhaustive baseline & All 1,100 forward DAGs with $N\le5$;
26,706 resource cases, plus 100 seeded fan-in-two random graphs.\\
Independent solver agreement & 17,378 instances: all 684 fan-in-two
DAGs with $N\le5$ at all tested resources, plus 1,000 seeded larger instances.\\
Static/operational bridge & 264,121 finite time assignments or replay cases
on all 75 DAGs with $1\le N\le4$.\\
Fixed-width threshold family & 216 horizon-$K$ decisions, $2\le K\le10$,
all windows, widths two and three; independent witness replay.\\
Depth-two threshold family & 25 programs, 5,720 horizon-two decisions,
$2\le m\le10$, all allowed $r$ and all windows/widths.\\
Approximate saturation family & 2,020 exact resource comparisons for
$1\le K\le10$, including oversized windows; 1,185 tolerance thresholds.\\
Compact certificates & 350 independent random instances and
24 deliberately malformed certificates.\\\bottomrule
\end{tabular}
\end{center}

Lean 4.34.1 (compiler commit
\texttt{5045d0056413266e57c625dcd7c365b10e377c52}) uses Std, with no Mathlib
dependency. The main modules expose 68 audited declarations; the resource-footprint
module adds 18, three supplementary Bellman/search modules add 20,
the conditional randomized-potential module adds 10, and the actual
repeated-gadget module adds 20, and work-count/corridor modules add 25
distinct declarations (161 in total). Root
recompilation uses fresh temporary compiled objects. Audits permit only
Lean's standard \texttt{propext}, \texttt{Classical.choice} and
\texttt{Quot.sound}; no new unproved axiom, \texttt{sorry}, or \texttt{admit}
is used. A documented process-path compatibility shim is a host runtime
accommodation and does not change the compiler or kernel.

The formal scope includes operational/time-assignment equivalence, finite
optimality, resource caps, the exact fixed-schedule minimum-window formula,
a concrete executable time-and-lane enumeration
algorithm with soundness/completeness/optimality, operational Bellman lower
certificates, a general attaining-potential theorem, and the arbitrary-local-
policy four-job obstruction. Generic finite-search lemmas and the actual
concrete enumerator are distinguished. These proofs do not verify the
Python program, the polynomial forced-segment DP, the corridor candidate
count or parameterized algorithm, the full general
online theorem, or the unbounded saturation families. Finite tests are not
universal proofs. Independent agent review is not academic peer review.

For the new polynomial optimizer, independent time-variable constraints
and ROB replay agree on 4,472 instances; 26,746 clear-prefix values also
agree. A separately implemented general-excess compressed algorithm was
checked on 13,380 cases. The corridor implementation has 4,797 independent
constraint comparisons, 1,921 feasible witness replays, and 9,384 boundary
checks on 2,050 complete legal time maps. The linear saturation theorem
has separate graph, abstract macro-pattern and negative-certificate audits.
Each saved report gives its actual ranges and seeds; these counts overlap.

\section{Limitations and future work}
The model omits variable operation latency, functional-unit classes, finite
retirement bandwidth, physical-register allocation, memory hierarchy,
speculation, branches, exceptions, dynamic power and clock constraints.
Its online information model is deliberately explicit and may differ from
what a concrete processor can infer. No implementation-level speed or energy
claim follows from these theorems.

The next questions are fixed-parameter tractability in $B-W$ alone,
efficient optimality certificates for the compressed graph, tight online ratios
at other fixed resources, composition of larger hidden-subset families, efficient
certificates at finite issue width, and formalization of the general
resource and observation theorems. A deep originality audit remains needed,
especially for bounded-lookahead and FIFO storage scheduling. All research
states and reproducible commands are saved in \texttt{STATUS.md} and
\texttt{HANDOFF.md}; the authorized research branch is
\texttt{partner/dist-ILP1-20261014}.

\begin{thebibliography}{10}
\small
\bibitem{wall} D. W. Wall. \emph{Limits of Instruction-Level Parallelism}.
WRL Research Report 93/6, 1993.
\url{https://davidwall.info/papers/WRL-TR-93.6.pdf}.
\bibitem{dubey} P. K. Dubey, G. B. Adams III, M. J. Flynn.
\emph{Instruction Window Size Trade-Offs and Characterization of Program
Parallelism}. IEEE Transactions on Computers, 1994.
\url{https://doi.org/10.1109/12.278481}.
Institutional record inspected; full-text access incomplete.
\bibitem{fields} B. Fields, S. Rubin, R. Bod\'ik.
\emph{Focusing Processor Policies via Critical-Path Prediction}. ISCA, 2001.
\url{https://www.eecg.utoronto.ca/~moshovos/ACA06/readings/critical.pdf}.
\bibitem{tullsen} D. M. Tullsen, B. Calder.
\emph{Computing Along the Critical Path}. UCSD Technical Report, October 1998.
\url{https://cseweb.ucsd.edu/~calder/papers/UCSD-CSE98-CPProfile.pdf}.
\bibitem{navada} S. Navada, N. K. Choudhary, E. Rotenberg.
\emph{Criticality-driven Superscalar Design Space Exploration}. PACT, 2010.
\url{https://ericrotenberg.wordpress.ncsu.edu/files/2025/06/paper_PACT-19.pdf}.
\bibitem{cg} E. G. Coffman Jr., R. L. Graham.
\emph{Optimal Scheduling for Two-Processor Systems}. Acta Informatica 1,
200--213, 1972. \url{https://doi.org/10.1007/BF00288685}.
\bibitem{cire} A. A. Cire, W.-J. van Hoeve.
\emph{Multivalued Decision Diagrams for Sequencing Problems}.
Author manuscript revised May2013, section9, Lemma8 and Theorem3.
\url{https://www.andrew.cmu.edu/user/vanhoeve/papers/mdd_sequencing.pdf}.
\bibitem{graham} R. L. Graham. \emph{Bounds for Certain Multiprocessing
Anomalies}. Bell System Technical Journal 45, 1563--1581, 1966.
\url{https://doi.org/10.1002/j.1538-7305.1966.tb01709.x}.
\bibitem{oai} OpenAI. \emph{Three-machine unit-time precedence scheduling},
family 124, September 2026 artifact. Inspected repository commit
\texttt{1581d4cbdd69a0faf34b2089614fe2a814316a6a}.
\url{https://github.com/openai/math/blob/1581d4cbdd69a0faf34b2089614fe2a814316a6a/lean/docs/124.md}.
\end{thebibliography}
\end{document}
